% Abstract text (150--250 words). Keywords are set in main.tex.
Classifications of cellular automata group rules by their behavior, but they do not measure how similar the rules are. We look for a similarity index based on the evolution that allows rules to be compared, even when they have different numbers of states or neighbors. First, we define four exact relations (state permutation, mirror, state reduction and neighborhood reduction), which reduce the 256 rules of the elementary cellular automata (ECA) to 81 classes. Then we compare the transition graph on a ring of $n$ cells: its canonical form, together with the quotient by translations, distinguishes the 88 classes under permutation and mirror with only 6 cells, but the comparison is binary. As a gradual measure, we propose spectral moments with attractor mass, a fixed-length vector that, unlike the eigenvalues and the spectral moments, counts how many configurations reach each attractor and in how many steps. With the cosine distance, this vector distinguishes 86 of the 88 classes (87 together with the quotient), including rules 0 and 8, and gives zero distance to the relations that preserve the graph and, in the quotient, to those that differ by a translation, but it does not identify state reduction. Moreover, it depends on the size of the space, and it groups the ECA into four families according to the dominant period of their attractors.
